\documentclass[10pt,a4paper]{amsart}
\usepackage[margin=19mm]{geometry}
\usepackage{amsmath,amssymb,mathtools}
\usepackage[scr=rsfs,cal=euler]{mathalfa}
\usepackage{xfrac}
\usepackage[unicode,hidelinks,bookmarks=false]{hyperref}
\newcommand{\ie}{i.e.\ }
\newcommand{\cf}{cf.\ }
\newcommand{\resp}{respectively\ }
\newcommand{\Subsection}{\subsection}
\providecommand{\nobreakdash}{\nobreak\hbox{-}}
\setlength{\emergencystretch}{2em}
\multlinegap0pt
\usepackage{amssymb}
\newtheorem{theorem}{Theorem}
\newcommand{\bona}{\mathrm{BoNa}}
\title{Boolean-Narayana numbers}
\author{Mikl\'os B\'ona}
\date{Source: arXiv:2602.11355v5 (CC BY 4.0)}
\begin{document}
\maketitle

\noindent\emph{Source and licence.} Boolean-Narayana numbers. Version arXiv:2602.11355v5 (CC BY 4.0), distributed by arXiv under the Creative Commons Attribution 4.0 International licence, http://creativecommons.org/licenses/by/4.0/, as recorded in the arXiv licence metadata for this exact version and shown on the abstract page https://arxiv.org/abs/2602.11355v5. Full licence notice: \texttt{licenses/arxiv-2602.11355-CC-BY-4.0.txt}.\par\smallskip

\noindent\textit{Editorial scope.} Counted unit: the article's theorem explicit (source0.tex lines 187--190). The article's complete original proof of it is delivered with it (source0.tex lines 192--217, one \texttt{proof} environment of 26 lines). No mathematics is written, repaired or re-derived: the statement and the proof are the article's own lines, in the article's own notation, and no retained line is substituted or re-flowed.  Retained uncounted as the source-local context the proof needs: lines 167--174.  Nothing was omitted from any retained span, so the package carries no omission record.  No retained line contains a citation, so the wrapper carries no bibliography and no bibliography entry.  No retained line contains a figure environment, an image-inclusion command or drawing code, so no image is delivered and the package declares no asset dependency.  Editorial formatting only, in addition: an AMS \texttt{amsart} preamble that restates the article's own declarations and macros used by the retained lines, namely the environments \texttt{theorem} on separate counters, together with the standard packages those lines need.  The retained environments print in this document as Theorem 1 in order of appearance, whereas the article prints the same items with the numbering of its own full document.  The complete native source of the article as distributed in the arXiv e-print for this exact version is bundled unchanged as \texttt{sources/194470/source0.tex}.
\begin{theorem} \label{lif}  Let 
$F(z)$ be the unique formal power series with
$F(0)=0$ satisfying 
\begin{equation} \label{lif-def} F(z)=z\phi(F(z),\end{equation}
where $\phi(y)$ is a formal power series with $\phi(0)\neq 0$. 
Then for all positive integers $n$, the equality
\[[z^n]F(z)=\frac{1}{n}[y^{n-1}\phi(y)^n] \] holds.
\end{theorem}

\begin{theorem}  \label{explicit}
Let $k\leq (n+1)/2$, where $n$ and $k$ are positive integers. Then
\[\bona(n,k)=\frac{1}{n} {n\choose k-1} \sum_{j=0}^{k-1} 2^j {k-1\choose j}{n-k+1\choose j+1} .\]
\end{theorem}

\begin{proof} Let $t(n,k)=\bona(n,k+1)$ be the number of 0-1 trees on $n$ vertices with $k$ right edges, and let
\begin{equation} \label{bivariate-tree} T(z,u)=\sum_{n\geq 1,k\geq 0} t(n,k) z^n u^k = \sum_{n\geq 1,k\geq 1} \bona(n,k)z^n u^{k-1}. \end{equation} 
Removing the root of a 0-1 tree we either get the empty set, or one tree that was connected to the root by a left edge, or one tree that was connected to the
root by a right edge, or two trees, one of which was connected to the root by a left edge.  This implies the functional equation
\[T=z(1+uT+T+2uT^2).\]
In order to use the Lagrange Inversion Formula, we set 
\begin{equation} \label{defofphi} \phi(y)=\phi(y,u)=1+(u+1)y+2uy^2,\end{equation} and then (\ref{defofphi}) becomes
\begin{equation} \label{lagrange-first} T=z\phi(T)) .\end{equation} 

By Theorem \ref{lif}, 
\begin{equation} \label{lagrange-second} [z^nu^{k-1}] T(z,u)=\frac{1}{n}[y^{n-1}u^{k-1}] \phi(y)^n.\end{equation}

Now notice that 
\[\phi(y)=(1+y)+u(y+2y^2)=(1+y)+uy(1+2y) .\]
Therefore, for $\phi(y)^n$ on the right-hand side of (\ref{lagrange-second}), we take the power $((1+y)+uy(1+2y))^n$, and notice
that the term $u^{k-1}$ will come from the summand of that binomial expansion in which the $u$-term is taken to the $(k-1)$st power. 
That is, 
\[[u^{k-1}]\phi(y)^n = {n\choose k-1} (1+y)^{n-k+1}y(1+2y))^{k-1} .\] 
Recalling (\ref{lagrange-second}), this means
\begin{eqnarray*} \bona(n,k) & = & [z^nu^{k-1}]T(z,u)=\frac{1}{n} {n\choose k-1}[y^{n-1} ]y^{k-1}\cdot(1+2y)^{k-1} (1+y)^{n-k+1} \\
& = & [z^nu^{k-1}]T(z,u)=\frac{1}{n} {n\choose k-1}[y^{n-k} ](1+2y)^{k-1} (1+y)^{n-k+1}.\end{eqnarray*}
Finally, expanding the binomial sums we get the expression
\[\bona(n,k)=\frac{1}{n} {n\choose k-1} \sum_{j=0}^{\min(k-1,n-k)} 2^j {k-1\choose j}{n-k+1\choose j+1} ,\]
However, recall that we assumed that $k\leq (n+1)/2$, (which, for symmetry reasons, does not result in loss of generality), so 
$k-1\leq n-k$, and our theorem is proved.  
\end{proof}

\end{document}
